\subsubsection{Desplazamiento virtual en péndulo sobre columna rotatoria}

Una masa $m$ (punto $P$) está suspendida mediante una varilla rígida de longitud $L_2$ desde el extremo de un brazo horizontal de longitud $L_1$. El brazo está unido rígidamente a una columna vertical que rota con una velocidad angular $\dot{\phi}$ alrededor del eje $Z$. El péndulo oscila en el plano vertical generado por la columna y el brazo, definido por el ángulo $\theta$ con la vertical. Considerando que la gravedad actúa hacia abajo, realice lo siguiente:

\begin{enumerate}
    \item Defina los vectores de posición del punto $P$ en un sistema de coordenadas inercial.
    \item Determine el vector de velocidad del punto $P$.
    \item Determine la expresión para el \textbf{desplazamiento virtual} $\delta \vec{r}_P$ del punto $P$ en términos de las coordenadas generalizadas y sus variaciones.
\end{enumerate}

\begin{figure}[htbp]
    \centering
    \begin{tikzpicture}[>=Latex, scale=1.4]
        % --- Parámetros ---
        \def\ColW{0.8}    % Ancho columna
        \def\ColH{3.5}    % Alto columna
        \def\Lone{2.8}    % Longitud L1
        \def\Ltwo{2.2}    % Longitud L2
        \def\angTheta{25} % Ángulo theta
        
        % Coordenadas
        \coordinate (Base) at (0,0);
        \coordinate (Top) at (0, \ColH);
        \coordinate (ArmStart) at (0, \ColH-0.5);
        \coordinate (ArmEnd) at (\Lone, \ColH-0.5);
        \coordinate (Mass) at ($(ArmEnd) + ({sin(\angTheta)*\Ltwo}, {-cos(\angTheta)*\Ltwo})$);

        % --- 1. Base y Columna ---
        % Base soporte
        \fill[pattern=north east lines] (-0.8, -0.2) rectangle (0.8, 0);
        \draw[thick] (-0.8, 0) -- (0.8, 0);
        %\draw[fill=gray!20] (0,0.1) ellipse (0.9 and 0.15); % Base de la columna visual

        % Columna (Cilindro)
        \fill[left color=gray!30, right color=gray!10, middle color=gray!5] (-\ColW/2,0) rectangle (\ColW/2, \ColH);
        \draw[thick] (-\ColW/2, 0) rectangle (\ColW/2, \ColH);
        
        % Tapa columna (Elipse superior)
        %\draw[fill=gray!20] (0, \ColH) ellipse (\ColW/2 and 0.15);
        
        % Eje Z
        \draw[->, thick] (0, \ColH) -- (0, \ColH+1.0) node[right] {$Z$};

        % --- 2. Brazo L1 ---
        \draw[line width=4pt, gray!40] (0, \ColH-0.6) -- (0, \ColH-0.4) -- (\Lone, \ColH-0.4) -- (\Lone, \ColH-0.6) -- cycle;
        \draw[thick] (0, \ColH-0.6) -- (\Lone, \ColH-0.6);
        \draw[thick] (0, \ColH-0.4) -- (\Lone, \ColH-0.4);
        \draw[thick] (\Lone, \ColH-0.6) -- (\Lone, \ColH-0.4);
        
        % Cota L1
        \draw[|<->|] (0, \ColH-0.2) -- (\Lone, \ColH-0.2) node[midway, above] {$L_1$};

        % --- 3. Péndulo L2 y Masa ---
        \draw[thick] (ArmEnd) -- (Mass) node[midway, right] {$L_2$};
        
        % Masa P
        \shadedraw[ball color=gray!30!white] (Mass) circle (0.25) node[] (Mass) {$m$};
        \node[left, xshift=-8pt] at (Mass) {$P$};

        % --- 4. Ángulos y Rotación ---
        % Referencia Vertical en extremo del brazo
        \draw[dashed] (ArmEnd) -- ($(ArmEnd) + (0, -2.0)$);
        
        % Ángulo Theta
        \draw[->] ($(ArmEnd) + (0, -1.2)$) arc (270:{270+\angTheta}:1.2);
        \node at ($(ArmEnd) + ({sin(\angTheta/2)*0.9}, {-cos(\angTheta/2)*0.9})$) [] {$\theta$};

        % Rotación Phi
        \draw[->, thick] (-0.8, 1.5) arc (180:360:0.8 and 0.3);
        \node[right] at (0.8, 1.5) {$\dot{\phi}$};

        % --- 5. Gravedad ---
        \draw[->, thick] (2.5, 2.5) -- (2.5, 1.5) node[midway, right] {$g$};

    \end{tikzpicture}
    \caption{Sistema de péndulo suspendido de un brazo en rotación acimutal.}
    \label{fig:pendulo_rotatorio}
\end{figure}